\documentclass[12pt,a4paper]{article}

% 這是一份latex 文件的 preamble
% 請AI閱讀後,將附上的PDF整理,以latex code整理,儘量用longtable
% 請注意我的名字: 謝慕揚, MD, PhD, FESC
% 表格請注意換行,不該在cell內使用 \\
% 表格請不要有垂直分隔線 |
% 所有 longtable 中的儲存格內換行都改用 \newline, 不能直接用\\
% 整理中若碰到藥物名稱,請使用英文,不要將藥物翻成中文。同樣,檢驗與檢查,請維持英文。


% Including essential packages for Chinese typesetting
\usepackage{xeCJK}
\usepackage{fontspec}
% Configuring Chinese font with Noto Serif CJK TC, fallback to SimSun
\IfFontExistsTF{Noto Serif CJK TC}{
  \setCJKmainfont{Noto Serif CJK TC}
  \setCJKsansfont{Noto Serif CJK TC}
  \setCJKmonofont{Noto Serif CJK TC}
}{\setCJKmainfont{SimSun}
  \setCJKsansfont{SimSun}
  \setCJKmonofont{SimSun}
}

% Including other necessary packages
\usepackage{geometry}
\usepackage{titlesec}
\usepackage{enumitem}
\usepackage{xcolor}
\usepackage{colortbl}
\usepackage{tcolorbox}
\usepackage{booktabs}
\usepackage{longtable}
\usepackage{array}
\usepackage{graphicx}
\usepackage{tikz}
\usepackage{fancyhdr}
\usepackage{multicol}
\usepackage{datetime2}
\usepackage{hyperref}
\usepackage{mdframed}
\usepackage{amsmath}

\hypersetup{
    colorlinks=true,
    linkcolor=emergency_blue,
    citecolor=emergency_blue,
    urlcolor=emergency_blue,
    pdfborder={0 0 0}
}

% Defining page geometry
\geometry{
  top=2.5cm,
  bottom=2.5cm,
  left=2cm,
  right=2cm,
  headheight=25pt
}

% Adjusting topmargin to compensate for increased headheight
\addtolength{\topmargin}{-10pt}

% Defining colors
\definecolor{emergency_red}{RGB}{186,24,27}
\definecolor{emergency_blue}{RGB}{0,114,188}
\definecolor{emergency_gray}{RGB}{120,120,120}
\definecolor{highlight_yellow}{RGB}{255,250,205}

% Setting title formats
\titleformat{\section}[block]{\Large\bfseries\color{emergency_red}}{\thesection}{10pt}{}
\titleformat{\subsection}[block]{\large\bfseries\color{emergency_blue}}{\thesubsection}{8pt}{}
\titleformat{\subsubsection}[block]{\normalsize\bfseries\color{emergency_gray}}{\thesubsubsection}{6pt}{}

% Defining custom environment for important notes
\newmdenv[
    linecolor=emergency_red,
    backgroundcolor=highlight_yellow,
    linewidth=2pt,
    topline=true,
    bottomline=true,
    leftline=true,
    rightline=true,
    innertopmargin=10pt,
    innerbottommargin=10pt,
    innerleftmargin=10pt,
    innerrightmargin=10pt
]{importantbox}

% Defining note command with tcolorbox
\newcommand{\note}[1]{
    \par
    \noindent
    \begin{tcolorbox}[
        colback=emergency_blue!5!white,
        colframe=emergency_blue,
        title=\textbf{重點提醒},
        fonttitle=\bfseries,
        boxrule=0.5pt,
        arc=3pt,
        left=6pt,
        right=6pt,
        top=6pt,
        bottom=6pt
    ]
    #1
    \end{tcolorbox}
}

% Setting up header and footer
\pagestyle{fancy}
\fancyhf{}
\fancyhead[L]{\textcolor{emergency_red}{新竹臺大分院內科部教學文件}}
\fancyhead[R]{\textcolor{emergency_blue}{住院醫師專業素養與工作流程}}
\fancyfoot[C]{\textcolor{emergency_gray}{\thepage}}
\renewcommand{\headrulewidth}{1pt}
\renewcommand{\headrule}{\hbox to\headwidth{\color{emergency_red}\leaders\hrule height \headrulewidth\hfill}}

% Defining custom date format for ISO 8601
\DTMnewdatestyle{iso}{
  \renewcommand{\DTMdisplaydate}[4]{\number##1-\number##2-\number##3}
  \renewcommand{\DTMDisplaydate}[4]{\number##1-\number##2-\number##3}
}
\DTMsetdatestyle{iso}
\newcommand{\mydate}{\DTMtoday}

\begin{document}

\begin{center}
{\Huge\textcolor{emergency_red}{\textbf{住院醫師專業素養}}} \\[10pt]
{\Huge\textcolor{emergency_red}{\textbf{與工作流程規範}}} \\[15pt]
{\Large\textcolor{emergency_blue}{新竹臺大分院內科部教學訓練}} \\[20pt]
{\large 教學指導：謝慕揚, MD, PhD, FESC} \\[5pt]
{\large 新竹臺大分院心血管中心} \\[5pt]
{\large 教學部副主任} \\[20pt]
{\normalsize 日期：\mydate}
\end{center}

\vspace{1cm}

\begin{importantbox}
\textbf{教學目標}

本教學講義旨在建立住院醫師正確的專業態度與工作流程認知,確保病人照護品質與醫療團隊合作效能。所有內科住院醫師均應熟悉並落實本講義所述之規範。
\end{importantbox}

\newpage

\section{住院醫師的專業素養核心概念}

\subsection{何謂專業素養(Professionalism)?}

專業素養是醫師職涯發展的基石,包含以下核心要素:

\begin{itemize}[leftmargin=*]
\item \textbf{病人優先(Patient-Centered Care):} 病人的安全與福祉永遠是第一優先
\item \textbf{責任感(Accountability):} 對自己負責的病人承擔完整照護責任
\item \textbf{主動學習(Active Learning):} 主動尋求學習機會,不僅被動等待指示
\item \textbf{團隊合作(Teamwork):} 與護理師、其他醫師密切合作
\item \textbf{溝通技巧(Communication):} 與病人、家屬、團隊成員保持良好溝通
\item \textbf{時間管理(Time Management):} 有效率地完成工作,不影響團隊運作
\item \textbf{終身學習(Lifelong Learning):} 持續精進醫學知識與臨床技能
\end{itemize}

\note{
\textbf{重要觀念:}住院醫師不是「醫學生」,而是「正在接受訓練的醫師」。你們已經是醫療團隊的正式成員,需要承擔相應的責任與義務。
}

\subsection{常見的專業素養問題}

以下是住院醫師訓練過程中常見的問題,請各位自我檢視:

\begin{itemize}[leftmargin=*]
\item 認為「下班時間」就可以不處理病人問題
\item 只待在護理站,很少實際去看病人
\item 等待主治醫師或資深醫師告知該做什麼,缺乏主動性
\item 不熟悉自己負責的病人狀況
\item 查房時站得很遠,不清楚主治醫師在詢問自己
\item 對病房系統操作流程不熟悉,但不主動學習
\end{itemize}

\begin{importantbox}
\textbf{態度決定高度}

醫學知識可以慢慢學習,但專業態度必須從第一天就建立。如果基本態度有問題,再多的知識也無法成為一位好醫師。
\end{importantbox}

\newpage

\section{住院醫師每日工作流程標準}

\subsection{接新病人的標準流程}

當你接到新病人時,必須完成以下步驟:

\begin{enumerate}[leftmargin=*]
\item \textbf{第一時間親自看病人}
\begin{itemize}
\item 不是只在急診或病房護理站聽交班
\item 必須親自到病人床邊進行身體檢查
\item 與病人及家屬建立初步關係
\item 確認病人的主訴與目前狀況
\end{itemize}

\item \textbf{仔細查看所有檢查報告}
\begin{itemize}
\item 不能只聽數值,要親自看影像與報告
\item Lab data、X-ray、CT、心電圖等都要自己判讀
\item 如有疑問,主動詢問資深醫師或放射科醫師
\end{itemize}

\item \textbf{建立完整病歷}
\begin{itemize}
\item History taking要詳細且有邏輯
\item Physical examination要確實執行並記錄
\item Assessment and Plan要有明確思路
\end{itemize}

\item \textbf{撰寫Admission Note}
\begin{itemize}
\item 必須在病人入院後24小時內完成
\item 格式要完整:Chief Complaint, Present Illness, Past History, Physical Examination, Lab Data, Assessment, Plan
\item 計畫要具體可執行,不能只寫「持續觀察」
\end{itemize}

\item \textbf{與主治醫師討論治療計畫}
\begin{itemize}
\item 主動向VS報告病人狀況
\item 提出自己的評估與計畫
\item 詢問不清楚的治療策略
\end{itemize}
\end{enumerate}

\note{
\textbf{關鍵概念:}接病人不是「交接資料」,而是「接手照護責任」。從你接手的那一刻起,這位病人的安危就與你密切相關。
}

\subsection{每日查房流程}

每天早上的查房是住院醫師最重要的工作之一:

\begin{enumerate}[leftmargin=*]
\item \textbf{查房前準備(晨會前)}
\begin{itemize}
\item 提早到病房,先看過夜班護理記錄
\item 確認過夜有無特殊事件
\item 查看當天的Lab results
\item 準備好要向VS報告的內容
\end{itemize}

\item \textbf{查房時的定位}
\begin{itemize}
\item 站在病床旁邊,靠近VS的位置
\item 保持警覺,注意VS何時要詢問你
\item 準備好病人的最新資訊
\item 不要只站在門口或走廊
\end{itemize}

\item \textbf{報告病人狀況}
\begin{itemize}
\item 簡潔扼要地說明病人昨天的變化
\item 重要的檢驗數據要能立即回答
\item 如果不清楚,誠實說不知道,但要主動查詢
\item 提出今天的處置計畫
\end{itemize}

\item \textbf{執行VS的醫囑}
\begin{itemize}
\item 查房後立即開立醫囑或安排檢查
\item 不能「等到下午有空再開」
\item 如有疑問,立即詢問VS或資深住院醫師
\end{itemize}

\item \textbf{查房後追蹤}
\begin{itemize}
\item 與護理師確認醫囑執行狀況
\item 追蹤檢查報告是否出來
\item 評估病人對治療的反應
\end{itemize}
\end{enumerate}

\begin{importantbox}
\textbf{常見錯誤:}「我以為VS不是在問我」

在教學醫院,VS提出問題時,預設就是在詢問負責該病人的住院醫師。即使你站得較遠,也應該保持注意力,隨時準備回應。
\end{importantbox}

\subsection{Progress Note 撰寫規範}

Progress Note是記錄病人每日變化與治療決策的重要文件,也是法律文書。

\subsubsection{Progress Note的基本格式}

標準的Progress Note應包含以下要素:

\begin{itemize}[leftmargin=*]
\item \textbf{S (Subjective):}病人的主觀症狀
\begin{itemize}
\item 病人今天的感覺如何?
\item 有無新的不適?
\item 對治療的反應?
\end{itemize}

\item \textbf{O (Objective):}客觀的檢查數據
\begin{itemize}
\item Vital signs (最新的體溫、血壓、心跳、呼吸、血氧)
\item Physical examination findings
\item Lab data (需要列出今天重要的檢驗結果)
\item Imaging studies results
\end{itemize}

\item \textbf{A (Assessment):}對病人狀況的評估
\begin{itemize}
\item 病人目前的問題列表
\item 每個問題的嚴重程度與趨勢
\item 治療效果的評估
\end{itemize}

\item \textbf{P (Plan):}今日的處置計畫
\begin{itemize}
\item 針對每個問題的具體處置
\item 需要安排的檢查
\item 藥物調整
\item 預計的處置時程
\end{itemize}
\end{itemize}

\subsubsection{Progress Note的撰寫時機}

\begin{itemize}[leftmargin=*]
\item \textbf{每日查房後:}必須在查房後2-4小時內完成當日的Progress Note
\item \textbf{病情變化時:}病人有重大變化時,應立即補充記錄
\item \textbf{特殊處置後:}執行invasive procedure後需記錄過程與結果
\item \textbf{會診後:}會診其他科別後應記錄會診建議與處置
\end{itemize}

\note{
\textbf{系統操作:}如果對內部Progress Note系統不熟悉,應主動詢問資深醫師或資訊室。不能因為「不會用系統」而不寫Progress Note。
}

\subsubsection{Progress Note範例}

以下是一個標準的Progress Note範例:

\begin{mdframed}[linecolor=emergency_gray,linewidth=1pt]
\textbf{2025-11-07, Hospital Day \#3}

\textbf{S:}
Patient reports improved chest pain compared to yesterday. No dyspnea at rest. Denies palpitation or dizziness.

\textbf{O:}
\begin{itemize}[nosep]
\item Vital signs: T 36.8°C, BP 128/76 mmHg, HR 72 bpm, RR 16/min, SpO2 98\% on room air
\item Physical examination: clear breath sounds, regular heart beat without murmur, no pedal edema
\item Lab: Troponin-I 0.08 ng/mL (↓ from 2.45), CK-MB 3.2 ng/mL (↓ from 28.7), Creatinine 1.1 mg/dL
\item ECG: sinus rhythm, ST segment normalized in leads V2-V4
\end{itemize}

\textbf{A:}
\begin{enumerate}[nosep]
\item Acute coronary syndrome, NSTEMI, post-PCI to LAD - improving
\item Hypertension - controlled
\item Type 2 diabetes mellitus - under control
\end{enumerate}

\textbf{P:}
\begin{enumerate}[nosep]
\item Continue DAPT (Aspirin 100mg + Ticagrelor 90mg bid)
\item Continue statin (Atorvastatin 80mg qd)
\item Schedule for cardiac rehabilitation
\item Plan discharge tomorrow if condition remains stable
\item Arrange outpatient follow-up in 1 week
\end{enumerate}
\end{mdframed}

\newpage

% PAP格式Progress Note範例 - Peptic Ulcer Bleeding
% 可直接插入主文件中使用

\subsection{PAP格式Progress Note範例}

PAP (Problem-Assessment-Plan) 是另一種常用的Progress Note格式,特別適合用於有多個問題的病人。以下是消化性潰瘍出血的範例:

\begin{mdframed}[linecolor=emergency_gray,linewidth=1pt]
\textbf{2025-11-08, Hospital Day \#2}

\textbf{Problem \#1: Upper gastrointestinal bleeding secondary to peptic ulcer}

\textit{Assessment:}
\begin{itemize}[nosep,leftmargin=15pt]
\item 68-year-old male with history of chronic NSAID use
\item Presented with melena for 3 days and one episode of coffee-ground vomitus
\item Initial Hb 8.2 g/dL (baseline 13.5 g/dL), dropped to 7.1 g/dL
\item Emergency EGD performed yesterday revealed a Forrest IIa ulcer (2.5 cm) at gastric antrum with visible vessel
\item Endoscopic hemostasis achieved with hemoclip placement and epinephrine injection
\item No active bleeding noted post-procedure
\item Today: vital signs stable (BP 118/72 mmHg, HR 78 bpm), no further melena
\item Post-transfusion Hb 9.8 g/dL (received 2 units PRBC)
\item Nasogastric tube drainage: minimal blood-tinged fluid, no fresh blood
\end{itemize}

\textit{Plan:}
\begin{itemize}[nosep,leftmargin=15pt]
\item Continue high-dose PPI: Pantoprazole 40 mg IV q12h for at least 72 hours
\item NPO for now, may start clear liquid diet if no rebleeding for 24 hours
\item Continue hemodynamic monitoring: vital signs q4h, watch for signs of rebleeding
\item Follow-up Hb/Hct this afternoon
\item Discontinue all NSAIDs permanently
\item Target Hb >8 g/dL, transfuse if Hb <7 g/dL or symptomatic
\item Plan for repeat EGD in 48-72 hours to assess healing
\item Consult GI team if any signs of rebleeding
\item H. pylori test pending (serology and biopsy), treat if positive
\end{itemize}

\vspace{0.5cm}

\textbf{Problem \#2: Hypertension - controlled}

\textit{Assessment:}
\begin{itemize}[nosep,leftmargin=15pt]
\item Currently on Amlodipine 5 mg daily
\item BP readings stable: 118-132/68-78 mmHg over past 24 hours
\item No hypertensive emergency
\end{itemize}

\textit{Plan:}
\begin{itemize}[nosep,leftmargin=15pt]
\item Continue current antihypertensive regimen
\item Monitor BP q4h during acute bleeding phase
\item Resume oral medication when diet resumed
\end{itemize}

\vspace{0.5cm}

\textbf{Problem \#3: Chronic kidney disease, stage 3a}

\textit{Assessment:}
\begin{itemize}[nosep,leftmargin=15pt]
\item Baseline Creatinine 1.4-1.6 mg/dL (eGFR 45-52 mL/min/1.73m²)
\item Current Creatinine 1.5 mg/dL (stable from admission)
\item Adequate urine output (>50 mL/hr)
\item No signs of acute kidney injury
\item Pre-renal azotemia risk due to bleeding minimized by fluid resuscitation
\end{itemize}

\textit{Plan:}
\begin{itemize}[nosep,leftmargin=15pt]
\item Continue IV hydration: Normal saline at 80 mL/hr
\item Monitor daily renal function
\item Avoid nephrotoxic agents
\item Adjust medication doses according to renal function
\item Ensure adequate intravascular volume
\end{itemize}

\vspace{0.5cm}

\textbf{Problem \#4: Anemia secondary to acute blood loss}

\textit{Assessment:}
\begin{itemize}[nosep,leftmargin=15pt]
\item Admission Hb 8.2 g/dL → dropped to 7.1 g/dL → post-transfusion 9.8 g/dL
\item MCV 88 fL (normocytic), suggesting acute blood loss rather than chronic anemia
\item No signs of ongoing hemolysis
\item Reticulocyte count pending
\item Clinically stable: no tachycardia, no orthostatic hypotension
\end{itemize}

\textit{Plan:}
\begin{itemize}[nosep,leftmargin=15pt]
\item Serial Hb/Hct monitoring: check this afternoon, then q12h
\item Transfusion threshold: Hb <7 g/dL or symptomatic (chest pain, dyspnea, dizziness)
\item Start oral iron supplementation when diet resumed: Ferrous sulfate 325 mg tid
\item Monitor for signs of active bleeding: tachycardia, hypotension, melena, hematemesis
\end{itemize}

\end{mdframed}

\note{
\textbf{PAP格式的優點:}
\begin{itemize}[nosep]
\item 清楚區分每個問題,避免混淆
\item 每個問題都有獨立的評估與計畫
\item 適合複雜病人或多重問題的情況
\item 便於團隊成員快速找到特定問題的處置
\item 有助於確保所有問題都被適當處理
\end{itemize}
}

\begin{importantbox}
\textbf{撰寫PAP格式Progress Note的關鍵要點:}
\begin{enumerate}
\item \textbf{Problem列表要完整:}按照重要性或急迫性排序
\item \textbf{Assessment要有邏輯:}包含相關的檢查數據、理學檢查發現、病程變化
\item \textbf{Plan要具體可執行:}避免模糊的描述如「繼續觀察」
\item \textbf{動態調整:}每天根據病人狀況更新Problem List
\item \textbf{跨問題思考:}注意不同問題間的相互影響
\end{enumerate}
\end{importantbox}
\section{病人照護責任與交班邏輯}

\subsection{Primary Care Physician的概念}

作為住院醫師,你是病人的Primary Care Physician (PCP),這代表:

\begin{itemize}[leftmargin=*]
\item \textbf{24小時責任制:}即使下班,你仍對病人負有責任
\item \textbf{全面照護:}不只處理主要疾病,所有問題都要關注
\item \textbf{協調角色:}需要協調各專科會診與護理照護
\item \textbf{溝通橋樑:}與病人、家屬、團隊成員保持溝通
\end{itemize}

\begin{importantbox}
\textbf{重要觀念:}「我已經下班了,不處理病人問題」這種想法是完全錯誤的。

在醫院體制下,我們採用team-based care,但這不代表下班就可以完全不管。正確的做法是:
\begin{enumerate}
\item 下班前確實交班給on-call醫師
\item 如果on-call醫師諮詢你的意見,你仍應提供協助
\item 隔天早上要詢問過夜病人的狀況
\item 如果是重症病人,即使下班也要保持關注
\end{enumerate}
\end{importantbox}

\subsection{交班的標準流程}

良好的交班是確保病人安全的關鍵:

\subsubsection{日常交班(晨會)}

每天早上的晨會是整個醫療團隊的資訊同步時間:

\begin{enumerate}[leftmargin=*]
\item \textbf{報告格式}
\begin{itemize}
\item 病人基本資料:床號、姓名、年齡、診斷
\item 入院日期與住院天數
\item 目前主要問題與治療計畫
\item 需要團隊協助或注意的事項
\end{itemize}

\item \textbf{報告重點}
\begin{itemize}
\item 簡潔扼要,不要冗長
\item 突顯重要問題與變化
\item 提出需要討論的議題
\item 說明今日預計處置
\end{itemize}

\item \textbf{聆聽其他病人的報告}
\begin{itemize}
\item 即使不是你的病人,也要專心聽
\item 了解整個病房的動態
\item 學習其他病人的處置邏輯
\item 準備好提供跨病人支援
\end{itemize}
\end{enumerate}

\subsubsection{下班交班(Hand-off)}

下班前的交班要確保on-call醫師能順利接手:

\begin{enumerate}[leftmargin=*]
\item \textbf{交班內容}
\begin{itemize}
\item 所有自己負責病人的最新狀況
\item 今天已執行的處置與結果
\item 尚未完成的工作項目
\item 可能在夜間需要注意的問題
\item 重要的聯絡方式(家屬電話等)
\end{itemize}

\item \textbf{交班時機}
\begin{itemize}
\item 不能在最後一分鐘才交班
\item 預留足夠時間讓on-call醫師提問
\item 確認對方了解重點
\item 提供自己的聯絡方式以備諮詢
\end{itemize}

\item \textbf{交班工具}
\begin{itemize}
\item 使用標準化的交班表格或系統
\item 不能只靠口頭交班
\item 重要事項要寫下來
\item 確保資訊完整且可追溯
\end{itemize}
\end{enumerate}

\note{
\textbf{交班不是卸責:}交班給on-call醫師不代表你就完全沒有責任了。如果on-call醫師有疑問打電話給你,你應該提供協助,而不是說「我已經下班了」。
}

\subsection{跨團隊合作}

病房工作需要與多個團隊成員合作:

\begin{itemize}[leftmargin=*]
\item \textbf{護理師}
\begin{itemize}
\item 尊重護理師的專業判斷
\item 主動詢問病人的護理狀況
\item 及時回應護理師的諮詢
\item 感謝護理師的辛勞
\end{itemize}

\item \textbf{其他住院醫師}
\begin{itemize}
\item 主動分享學習心得
\item 互相支援臨床工作
\item 不要只關心自己的病人
\item 建立良好的同儕關係
\end{itemize}

\item \textbf{主治醫師}
\begin{itemize}
\item 主動報告病人狀況
\item 虛心接受指導與建議
\item 有問題立即詢問
\item 不要隱瞞錯誤或疏失
\end{itemize}

\item \textbf{其他專科}
\begin{itemize}
\item 會診時提供完整資訊
\item 尊重專科醫師的建議
\item 追蹤會診結果
\item 必要時進行跨科討論
\end{itemize}
\end{itemize}

\newpage

\section{內科住院醫師的核心技能}

\subsection{病史詢問(History Taking)}

完整且有邏輯的病史詢問是診斷的基礎:

\begin{enumerate}[leftmargin=*]
\item \textbf{Chief Complaint (CC)}
\begin{itemize}
\item 病人用自己的話描述主要症狀
\item 記錄症狀的持續時間
\item 一般只有1-2個主要症狀
\end{itemize}

\item \textbf{Present Illness (PI)}
\begin{itemize}
\item 症狀發生的時間序列
\item 每個症狀的特性(OPQRST)
\item Onset, Palliative/Provocative factors, Quality, Region/Radiation, Severity, Timing
\item 已接受的檢查與治療
\item 相關的positive與negative findings
\end{itemize}

\item \textbf{Past History}
\begin{itemize}
\item Past Medical History:重要的慢性疾病
\item Past Surgical History:過去的手術史
\item Medication History:目前使用的藥物
\item Allergy History:藥物或食物過敏
\end{itemize}

\item \textbf{Family History}
\begin{itemize}
\item 家族中的重要疾病
\item 特別注意遺傳性疾病
\end{itemize}

\item \textbf{Personal and Social History}
\begin{itemize}
\item 職業與工作環境
\item 抽菸、喝酒、檳榔等習慣
\item 居住環境與生活型態
\end{itemize}
\end{enumerate}

\subsection{身體檢查(Physical Examination)}

身體檢查要系統性且確實執行:

\begin{enumerate}[leftmargin=*]
\item \textbf{Vital Signs}
\begin{itemize}
\item 體溫(Temperature)
\item 血壓(Blood Pressure)
\item 心跳(Heart Rate)
\item 呼吸速率(Respiratory Rate)
\item 血氧飽和度(SpO2)
\end{itemize}

\item \textbf{General Appearance}
\begin{itemize}
\item 意識狀態
\item 營養狀態
\item 是否有急性病容
\end{itemize}

\item \textbf{HEENT (Head, Eyes, Ears, Nose, Throat)}
\begin{itemize}
\item 結膜蒼白或黃疸
\item 頸靜脈充盈
\item 甲狀腺腫大
\end{itemize}

\item \textbf{Chest}
\begin{itemize}
\item 呼吸音(breath sounds)
\item 心音(heart sounds)
\item 是否有雜音(murmur)或額外心音
\end{itemize}

\item \textbf{Abdomen}
\begin{itemize}
\item 視診、聽診、觸診、叩診
\item 注意壓痛點與反彈痛
\item 肝脾大小
\item 腸音
\end{itemize}

\item \textbf{Extremities}
\begin{itemize}
\item 水腫(edema)
\item 周邊脈搏
\item 皮膚溫度與顏色
\item 關節活動度
\end{itemize}

\item \textbf{Neurological Examination}(必要時)
\begin{itemize}
\item 意識程度與定向感
\item 肌力與肌張力
\item 深部肌腱反射
\item 感覺與運動功能
\end{itemize}
\end{enumerate}

\note{
\textbf{重要提醒:}身體檢查不能只寫「unremarkable」或「within normal limits」。必須寫出你確實執行了哪些檢查,以及具體的發現。
}

\subsection{臨床推理(Clinical Reasoning)}

好的臨床推理能力是成為優秀醫師的關鍵:

\begin{enumerate}[leftmargin=*]
\item \textbf{建立Problem List}
\begin{itemize}
\item 列出病人所有的問題
\item 區分急性與慢性問題
\item 判斷問題的優先順序
\end{itemize}

\item \textbf{建立Differential Diagnosis}
\begin{itemize}
\item 對每個問題列出可能的診斷
\item 按照可能性排序
\item 考慮life-threatening的診斷
\end{itemize}

\item \textbf{安排檢查}
\begin{itemize}
\item 基於診斷假設安排適當檢查
\item 不是「routine全套」檢查
\item 考慮檢查的效益與風險
\end{itemize}

\item \textbf{判讀檢查結果}
\begin{itemize}
\item 親自判讀,不只看報告
\item 結合臨床表現解讀數據
\item 注意趨勢變化,不只看單次數值
\end{itemize}

\item \textbf{制定治療計畫}
\begin{itemize}
\item 基於實證醫學的治療
\item 考慮病人的個別情況
\item 設定明確的治療目標
\item 規劃追蹤評估方式
\end{itemize}
\end{enumerate}

\newpage

\section{常見問題與解答(Q\&A)}

\subsection{關於工作時間與責任}

\begin{itemize}[leftmargin=*]
\item \textbf{Q1: 我下班後還需要管病人嗎?}

A1: 在team-based照護模式下,下班後由on-call醫師負責,但你仍需要:
\begin{enumerate}
\item 確實交班給on-call醫師
\item 如果on-call醫師有疑問,應提供協助
\item 隔天上班後主動詢問病人過夜狀況
\item 重症病人即使下班也要保持關注
\end{enumerate}

\item \textbf{Q2: 我不知道該做什麼,應該怎麼辦?}

A2: 主動詢問!不知道不是問題,不問才是問題。你可以:
\begin{enumerate}
\item 詢問VS或主治醫師
\item 請教資深住院醫師
\item 查閱教科書或clinical guidelines
\item 參加教學活動與討論
\end{enumerate}

\item \textbf{Q3: 我的時間不夠,無法完成所有工作怎麼辦?}

A3: 時間管理是需要學習的技能:
\begin{enumerate}
\item 早上提早到病房準備
\item 查房後立即處理重要事項
\item 善用零碎時間
\item 必要時請求同事支援
\item 與VS討論優先順序
\end{enumerate}
\end{itemize}

\subsection{關於系統操作}

\begin{itemize}[leftmargin=*]
\item \textbf{Q4: Progress Note系統不會用怎麼辦?}

A4: 立即學習!可以:
\begin{enumerate}
\item 請資深醫師或護理師示範
\item 查閱系統操作手冊
\item 聯絡資訊室求助
\item 參加系統教育訓練
\end{enumerate}
不能因為「不會用系統」而不寫Progress Note。

\item \textbf{Q5: 開立醫囑時不確定是否正確怎麼辦?}

A5: 病人安全第一!
\begin{enumerate}
\item 開立前再次確認
\item 不確定時詢問VS或資深醫師
\item 查閱藥物手冊或系統資訊
\item 與藥師討論
\item 寧可多問,不要猜測
\end{enumerate}
\end{itemize}

\subsection{關於學習與進步}

\begin{itemize}[leftmargin=*]
\item \textbf{Q6: 我覺得自己進步很慢怎麼辦?}

A6: 進步需要時間,但態度可以立刻改變:
\begin{enumerate}
\item 每天設定小目標
\item 主動尋求feedback
\item 從錯誤中學習
\item 參考優秀同儕的做法
\item 保持積極學習態度
\end{enumerate}

\item \textbf{Q7: 如何提升臨床能力?}

A7: 系統性的學習方式:
\begin{enumerate}
\item 參加晨會與教學活動
\item 認真準備每次報告
\item 閱讀相關文獻與指引
\item 參加臨床技能訓練
\item 請教資深醫師經驗
\item 反思每日的臨床決策
\end{enumerate}
\end{itemize}

\newpage

\section{自我檢核表}

請誠實評估自己在以下各項的表現,並設定改善目標:

\begin{longtable}{p{0.6\textwidth}p{0.15\textwidth}p{0.15\textwidth}}
\toprule
\textbf{項目} & \textbf{目前狀況} & \textbf{目標達成} \\
\midrule
\endfirsthead
\toprule
\textbf{項目} & \textbf{目前狀況} & \textbf{目標達成} \\
\midrule
\endhead
\midrule
\multicolumn{3}{r}{\textit{續下頁}} \\
\endfoot
\bottomrule
\endlastfoot

每天主動到病床看病人 & $\square$ 是 $\square$ 否 & $\square$ \\
熟悉所有負責病人的狀況 & $\square$ 是 $\square$ 否 & $\square$ \\
查房時專心聽講並積極回應 & $\square$ 是 $\square$ 否 & $\square$ \\
當日完成Progress Note撰寫 & $\square$ 是 $\square$ 否 & $\square$ \\
熟悉內部系統操作流程 & $\square$ 是 $\square$ 否 & $\square$ \\
主動詢問不懂的問題 & $\square$ 是 $\square$ 否 & $\square$ \\
確實執行交班流程 & $\square$ 是 $\square$ 否 & $\square$ \\
與護理師保持良好溝通 & $\square$ 是 $\square$ 否 & $\square$ \\
對病人照護保持責任感 & $\square$ 是 $\square$ 否 & $\square$ \\
準時參加晨會與教學活動 & $\square$ 是 $\square$ 否 & $\square$ \\
查房後立即開立醫囑 & $\square$ 是 $\square$ 否 & $\square$ \\
追蹤檢查報告並判讀 & $\square$ 是 $\square$ 否 & $\square$ \\
主動協助其他團隊成員 & $\square$ 是 $\square$ 否 & $\square$ \\
接受建議並改進缺點 & $\square$ 是 $\square$ 否 & $\square$ \\
保持積極學習態度 & $\square$ 是 $\square$ 否 & $\square$ \\
\end{longtable}

\section{結語與期許}

住院醫師訓練是醫師生涯中最辛苦但也最重要的階段。這段時間建立的專業素養與工作習慣,將影響你一輩子的職業生涯。

\begin{importantbox}
\textbf{三個核心價值}
\begin{enumerate}
\item \textbf{病人優先:}永遠把病人的安全與福祉放在第一位
\item \textbf{主動學習:}不要等待別人告訴你該做什麼,主動尋求學習機會
\item \textbf{團隊合作:}醫療是團隊工作,尊重每一位團隊成員
\end{enumerate}
\end{importantbox}

記住,成為一位好醫師的關鍵不只是醫學知識,更重要的是:

\begin{itemize}[leftmargin=*]
\item 對病人的同理心與關懷
\item 對醫療工作的熱忱與責任感
\item 終身學習的態度與習慣
\item 與團隊合作的能力
\item 面對挑戰的韌性與毅力
\end{itemize}


\vspace{2cm}

\begin{center}
\textcolor{emergency_blue}{--- 教學指導:謝慕揚 ---}\\[5pt]
\textcolor{emergency_gray}{新竹臺大分院心血管中心}\\[5pt]
\textcolor{emergency_gray}{教學部副主任}\\[5pt]
\textcolor{emergency_gray}{MD, PhD, FESC}
\end{center}

\end{document}